% GENERATED FILE. Edit canonical lessons, then run npm run generate:books.
% canonical-source-hash: fnv1a64:ac32661f278c4ced
% canonical-lessons: RU-C420-dosmotr, RU-C420-dosmatrivat, RU-C420-pogranichnik, RU-C420-deklaratsiya, RU-C420-deklarirovat

\chapter{Security and customs}
\label{ch:ru-security-and-customs}
\hlchaptermodality{420}

\begin{chapteropening}
\textbf{By the end of this chapter:} \emph{I can go through security and the border: the security check, searching bags, the border guard and the declaration.}

Complete the last lesson of chapter 420: I can go through security and the border: the security check, searching bags, the border guard and the declaration.
\end{chapteropening}

\section[\texorpdfstring{\ru{досмотр} (dosmótr)}{досмотр (dosmótr)}]{\ru{досмотр} (dosmótr) — a security check, an inspection}

\label{lesson:RU-C420-dosmotr}



\begin{quote}
Before the new one: say the Russian for luggage, then the Russian for boarding.
\end{quote}

\subsection*{You'll want to know: \ru{досмотр}}
\textbf{\ru{досмотр}} (\emph{dosmótr}) — \textquotedblleft{}a security check, an inspection\textquotedblright{}.

\begin{quote}
\ru{Перед} \ru{посадкой} \ru{все} \ru{проходят} \ru{досмотр}.
\end{quote}

\begin{quote}
\emph{Péred posádkoy vse prokhódyat dosmótr.}
\end{quote}

\begin{quote}
Before boarding everyone goes through the security check.
\end{quote}

Masculine. \emph{\ru{Таможенный} \ru{досмотр}} — a customs inspection.

\begin{cousinweb}
\emph{\ru{досмотр}} is built from \emph{\ru{до}-} 'thoroughly, to the end' and \emph{\ru{смотр}} 'a review', on \emph{\ru{смотреть}} 'to look', which EtymDB relates to Proto-Slavic \emph{motriti} 'to watch'.
\end{cousinweb}

\begin{grammarlens}[title={What you've built}]
One more word for telling, explaining and describing at B1.
\end{grammarlens}

\subsection*{Your turn}
\begin{itemize}
\raggedright
  \item \emph{Say it:} \emph{dosmótr}
  \item \emph{Say it:} \emph{dosmótr}, once more
  \item \emph{Say it:} a sentence of your own with \emph{dosmótr}
  \item \emph{Recall:} say the Russian for luggage, then the Russian for boarding, then say \emph{dosmótr} again
\end{itemize}

\begin{tcolorbox}[breakable,colback=teal!4,colframe=teal!35!black,title={Before you move on}]
What does \ru{досмотр} mean? (\textquotedblleft{}A security check, an inspection\textquotedblright{}.) Say it once more.
\end{tcolorbox}
\section[\texorpdfstring{\ru{досматривать} (dosmátrivat)}{досматривать (dosmátrivat)}]{\ru{досматривать} (dosmátrivat) — to search, to inspect (luggage)}

\label{lesson:RU-C420-dosmatrivat}



\begin{quote}
Before the new one: say the Russian for boarding, then the Russian for a security check.
\end{quote}

\subsection*{You'll want to know: \ru{досматривать}}
\textbf{\ru{досматривать}} (\emph{dosmátrivat}) — \textquotedblleft{}to search, to inspect (luggage)\textquotedblright{}.

\begin{quote}
\ru{На} \ru{границе} \ru{досматривают} \ru{все} \ru{сумки}.
\end{quote}

\begin{quote}
\emph{Na granítse dosmátrivayut vse súmki.}
\end{quote}

\begin{quote}
At the border they search all the bags.
\end{quote}

Imperfective; perfective \emph{\ru{досмотреть}}. Also to watch to the end: \emph{\ru{досмотреть} \ru{фильм}}.

\begin{cousinweb}
\emph{\ru{досматривать}} is built from \emph{\ru{до}-} 'thoroughly, to the end' and \emph{-\ru{сматривать}}, from \emph{\ru{смотреть}} 'to look', which EtymDB relates to Proto-Slavic \emph{motriti} 'to watch'.
\end{cousinweb}

\begin{grammarlens}[title={What you've built}]
One more word for telling, explaining and describing at B1.
\end{grammarlens}

\subsection*{Your turn}
\begin{itemize}
\raggedright
  \item \emph{Say it:} \emph{dosmátrivat}
  \item \emph{Say it:} \emph{dosmátrivat}, once more
  \item \emph{Say it:} a sentence of your own with \emph{dosmátrivat}
  \item \emph{Recall:} say the Russian for boarding, then the Russian for a security check, then say \emph{dosmátrivat} again
\end{itemize}

\begin{tcolorbox}[breakable,colback=teal!4,colframe=teal!35!black,title={Before you move on}]
What does \ru{досматривать} mean? (\textquotedblleft{}To search, to inspect (luggage)\textquotedblright{}.) Say it once more.
\end{tcolorbox}
\section[\texorpdfstring{\ru{пограничник} (pograníchnik)}{пограничник (pograníchnik)}]{\ru{пограничник} (pograníchnik) — a border guard}

\label{lesson:RU-C420-pogranichnik}



\begin{quote}
Before the new one: say the Russian for a security check, then the Russian for to search.
\end{quote}

\subsection*{You'll want to know: \ru{пограничник}}
\textbf{\ru{пограничник}} (\emph{pograníchnik}) — \textquotedblleft{}a border guard\textquotedblright{}.

\begin{quote}
\ru{Пограничник} \ru{долго} \ru{проверял} \ru{мой} \ru{паспорт}.
\end{quote}

\begin{quote}
\emph{Pograníchnik dólgo proveryál moy pásport.}
\end{quote}

\begin{quote}
The border guard checked my passport for a long time.
\end{quote}

Masculine. Built on \emph{\ru{граница}} — a border.

\begin{cousinweb}
\emph{\ru{пограничник}} is built on \emph{\ru{пограничный}} 'of the border', from \emph{\ru{граница}} 'a border', which is inherited from Proto-Slavic \emph{granica} (EtymDB).
\end{cousinweb}

\begin{grammarlens}[title={What you've built}]
One more word for telling, explaining and describing at B1.
\end{grammarlens}

\subsection*{Your turn}
\begin{itemize}
\raggedright
  \item \emph{Say it:} \emph{pograníchnik}
  \item \emph{Say it:} \emph{pograníchnik}, once more
  \item \emph{Say it:} a sentence of your own with \emph{pograníchnik}
  \item \emph{Recall:} say the Russian for a security check, then the Russian for to search, then say \emph{pograníchnik} again
\end{itemize}

\begin{tcolorbox}[breakable,colback=teal!4,colframe=teal!35!black,title={Before you move on}]
What does \ru{пограничник} mean? (\textquotedblleft{}A border guard\textquotedblright{}.) Say it once more.
\end{tcolorbox}
\section[\texorpdfstring{\ru{декларация} (deklarátsiya)}{декларация (deklarátsiya)}]{\ru{декларация} (deklarátsiya) — a customs declaration; a declaration}

\label{lesson:RU-C420-deklaratsiya}



\begin{quote}
Before the new one: say the Russian for to search, then the Russian for a border guard.
\end{quote}

\subsection*{You'll want to know: \ru{декларация}}
\textbf{\ru{декларация}} (\emph{deklarátsiya}) — \textquotedblleft{}a customs declaration; a declaration\textquotedblright{}.

\begin{quote}
\ru{В} \ru{самолёте} \ru{нам} \ru{дали} \ru{декларацию}.
\end{quote}

\begin{quote}
\emph{V samolyóte nam dáli deklarátsiyu.}
\end{quote}

\begin{quote}
On the plane they gave us a declaration form.
\end{quote}

Feminine. \emph{\ru{Заполнить} \ru{декларацию}} — to fill in a declaration.

\begin{cousinweb}
\emph{\ru{декларация}} is usually traced through French \emph{déclaration} to Latin \emph{dēclārātiō} 'a making clear' (EtymDB lists the French and the Latin words), on \emph{clārus} 'clear'.
\end{cousinweb}

\begin{grammarlens}[title={What you've built}]
One more word for telling, explaining and describing at B1.
\end{grammarlens}

\subsection*{Your turn}
\begin{itemize}
\raggedright
  \item \emph{Say it:} \emph{deklarátsiya}
  \item \emph{Say it:} \emph{deklarátsiya}, once more
  \item \emph{Say it:} a sentence of your own with \emph{deklarátsiya}
  \item \emph{Recall:} say the Russian for to search, then the Russian for a border guard, then say \emph{deklarátsiya} again
\end{itemize}

\begin{tcolorbox}[breakable,colback=teal!4,colframe=teal!35!black,title={Before you move on}]
What does \ru{декларация} mean? (\textquotedblleft{}A customs declaration; a declaration\textquotedblright{}.) Say it once more.
\end{tcolorbox}
\section[\texorpdfstring{\ru{декларировать} (deklarírovat)}{декларировать (deklarírovat)}]{\ru{декларировать} (deklarírovat) — to declare (at customs)}

\label{lesson:RU-C420-deklarirovat}



\begin{quote}
Before the new one: say the Russian for a border guard, then the Russian for a customs declaration.
\end{quote}

\subsection*{You'll want to know: \ru{декларировать}}
\textbf{\ru{декларировать}} (\emph{deklarírovat}) — \textquotedblleft{}to declare (at customs)\textquotedblright{}.

\begin{quote}
\ru{Деньги} \ru{выше} \ru{лимита} \ru{надо} \ru{декларировать}.
\end{quote}

\begin{quote}
\emph{Déngi výshe límita nádo deklarírovat.}
\end{quote}

\begin{quote}
Money above the limit has to be declared.
\end{quote}

The same form serves both aspects. Present \emph{\ru{декларирую}}.

\begin{cousinweb}
\emph{\ru{декларировать}} is usually traced to Latin \emph{dēclārāre} 'to make clear', on \emph{clārus} 'clear' (Lewis \& Short). EtymDB lists the Russian verb without a source.
\end{cousinweb}

\begin{grammarlens}[title={What you've built}]
One more word for telling, explaining and describing at B1.
\end{grammarlens}

\subsection*{Your turn}
\begin{itemize}
\raggedright
  \item \emph{Say it:} \emph{deklarírovat}
  \item \emph{Say it:} \emph{deklarírovat}, once more
  \item \emph{Say it:} a sentence of your own with \emph{deklarírovat}
  \item \emph{Recall:} say the Russian for a border guard, then the Russian for a customs declaration, then say \emph{deklarírovat} again
\end{itemize}

\begin{tcolorbox}[breakable,colback=teal!4,colframe=teal!35!black,title={Before you move on}]
What does \ru{декларировать} mean? (\textquotedblleft{}To declare (at customs)\textquotedblright{}.) Say it once more.
\end{tcolorbox}
